\section{Cross-Modality Cueing and Supervision}

The integration of radar and camera sensors has evolved beyond simple decision-level or feature-level fusion. Recent research explores **cross-modality supervision** and **transfer learning**, where one sensor modality assists the training or inference process of the other. This paradigm exploits the complementary strengths of the sensors: the semantic richness and high angular resolution of cameras, and the depth precision, velocity information, and weather robustness of radar. This section reviews methodologies focusing on cross-modal distillation, pseudo-labeling, and attention cueing.

\subsection{Radar-Supervised Visual Learning}
A primary challenge in monocular computer vision is the lack of absolute metric scale, particularly for depth estimation and scene flow. Radar, providing precise range and Doppler velocity measurements, serves as an effective source of weak supervision for camera-based networks.

Several approaches utilize sparse radar returns to supervise monocular depth estimation networks. Unlike LiDAR, radar data is sparse and noisy, requiring specific filtering mechanisms. \cite{Lin2020} proposed a self-supervised framework where radar depth points constitute a sparse ground truth signal, enforcing scale consistency on the dense depth map generated from monocular images. This allows the visual model to learn metric depth without requiring dense ground truth labels or stereo pairs.

Furthermore, the Doppler information from radar has been leveraged to supervise motion estimation. By projecting radar radial velocities onto the image plane, networks can act as a cross-modal teacher to constrain optical flow predictions, ensuring that the estimated visual motion aligns with the physical measurements from the radar \cite{Alonso2021}.

\subsection{Cross-Modal Knowledge Distillation}
Knowledge distillation (KD) transfers learned representations from a "teacher" modality (usually the one with higher information density or better labels) to a "student" modality.

\subsubsection{Camera-to-Radar Distillation}
Radar data is notoriously difficult for humans to annotate due to its non-intuitive representation (e.g., range-Doppler maps or sparse point clouds). Conversely, camera images are easily interpretable and have mature labeling ecosystems. Researchers have employed Cross-Modal Knowledge Distillation (CMKD) to transfer semantic knowledge from a camera network to a radar network.
\cite{Saha2021} demonstrated a method where a pre-trained image segmentation network supervises a radar detection network. By aligning the feature maps of the two modalities, the radar network learns to infer semantic classes (e.g., distinguishing pedestrians from static poles) even when the camera is not available at inference time or is degraded by darkness.

\subsubsection{Robustness Transfer for Adverse Weather}
A critical benefit of cross-modality learning is robustness transfer. While cameras degrade in fog, rain, or snow, radar remains stable. However, radar lacks the resolution to classify objects finely. \cite{Bijelic2020} introduced a deep sensor fusion architecture that uses multimodal data to learn robust representations. While primarily a fusion paper, it employs a mechanism where the network learns to rely on radar features when optical features are corrupted, effectively using radar to "supervise" the attention of the network in adverse weather conditions.
More explicitly, \cite{Gasperini2021} utilized cross-modal distillation to train a student network that operates on radar (or robust features) by mimicking a teacher network having access to clear-weather camera images, thereby extending performance into adverse scenarios.

\subsection{Cross-Modal Cueing and Attention}
In this methodology, one modality directs the computational focus of the other, often to improve efficiency or reduce false positives.

	extbf{Radar-Guided Visual Attention:} Radar detections can generate Region of Interest (RoI) proposals for image processing. \cite{Nabati2021} proposed CenterFusion, which, while a fusion method, utilizes a center-based radar association step. The radar detections are projected into the image frustum to cue the visual features, associating depth and velocity with visual objects. This "cueing" mechanism allows the visual detector to resolve depth ambiguities and filter out background visual clutter that lacks a corresponding radar return.

	extbf{Benefits of Cross-Modality Supervision:}
\begin{itemize}
    \item 	extbf{Label Efficiency:} Reduces the need for expensive manual annotation of radar data by leveraging auto-generated labels from camera-based foundation models.
    \item 	extbf{Metric Scale Recovery:} Enforces physical scale in vision networks using radar range measurements.
    \item 	extbf{Adverse Weather Robustness:} Facilitates the training of radar-based models that achieve semantic performance closer to cameras, maintaining functionality in fog and rain where cameras fail.
\end{itemize}
